\scanpage{9}
qu'on trouve l'anneau \inserted{cosimplicial} $Q^*$, et $Q^*[M]=Q^*\otimes_{\underline O_X}M$. Comme $X^{\widehat{[1]}}$ a une section sur $X^{\widehat{[0]}}$, on obtient aussitôt que \uline{$Q^*[M]$ est une résolution de} $M$ [en mettant sur ts les $Q^i[M]$ les structures extrême-droites de $\underline O_X$-Modules]. Il faut cependant préciser le sens qu'on attache à cette expression [les A.R.-pro-modules forment-ils une catégorie abélienne ?....] et s'assurer s'il ne faut pas une hypothèse de quasi-cohérence sur $M$.
En tous cas, \deleted{\unsure{pour}} \unsure{dénotant|désignant} par un $\ell$ le passage aux $\varprojlim$ des pro-objets, on trouve que $\ell Q^*[M]$ est une résolution de $M$ si $M$ \unsure{qu coh.}
\margin{8}
